\section{Empirical Noise-Based Significance of Distances}

    \subsection{Goal and Motivation}

        The goal of this procedure is to determine how likely a measured Euclidean distance between two gene expression vectors is explainable by noise alone. Here, ``noise'' refers to the combined effect of biological inter-individual variability and technical measurement uncertainty present in biological replicates.

        This is particularly important when integrating data from multiple sequencing projects, where different tissues or conditions (axes) often have very different numbers of biological replicates. For example, in cancer studies, tissues such as brain are typically much less sampled than tissues such as muscle. Without explicit noise modeling, such replicate imbalances lead to biased distance estimates and unreliable geometric interpretations.

        Our approach preserves the original semantic expression space while attaching an empirical confidence estimate to every measured distance.

    \subsection{Overview}

        All steps described below are performed independently for each axis (tissue or condition). No centering or rescaling of the expression vectors is applied. Instead, we estimate how much apparent distance can be generated by replicate variability alone.

        Let
        \[
            x_{g,a,i}
        \]
        denote the expression of gene $g$ on axis $a$ in biological replicate $i$.

        Let
        \[
            \bar{x}_{g,a} = \frac{1}{n_{g,a}} \sum_i x_{g,a,i}
        \]
        be the mean expression of gene $g$ on axis $a$ across its $n_{g,a}$ replicates.

    \subsection{Empirical Noise Model per Axis}

        For each axis $a$, we construct an empirical noise model as follows:

        \begin{enumerate}
            \item For every gene $g$, compute its mean expression $\bar{x}_{g,a}$.
            \item Collect all biological replicate values $x_{g,a,i}$ together with their mean expression values $\bar{x}_{g,a}$.
            \item This defines a scatterplot where the x-axis is mean gene expression $\bar{x}_{g,a}$ and the y-axis is biological replicates mapped to those mean expression levels $x_{g,a,i}$.
        \end{enumerate}

        This scatterplot encodes how biological replicate variability depends on mean expression level for \emph{that axis} (tissue).

    \subsection{Local k-Nearest-Neighbor Noise Sampling}

        For a given gene $g$ and axis $a$, we estimate noise at its mean expression level using adaptive k-nearest-neighbor (kNN) binning in mean-expression space. Note that Tensor Omics uses a k-d-tree of the mean expression levels to speed up lookup here.

        Let $k$ be a user-defined parameter (default $k=100$). This value is large enough to provide a stable empirical noise distribution while remaining local in expression space.

        For gene $g$ on axis $a$:

        \begin{enumerate}
            \item Find the $k$ nearest mean expression values $\bar{x}{g',a}$ to $\bar{x}{g,a}$ based on absolute difference in mean expression.
            \item Collect all biological replicate values $x_{g',a,i}$ belonging to those $k$ gene mean expression values.
            \item Convert these into \emph{signed} noise residuals:
                  \[
                      \epsilon = x_{g',a,i} - \bar{x}_{g',a}
                  \]
        \end{enumerate}

        The resulting set of $\epsilon$ values is the empirical noise distribution for expression level $\bar{x}_{g,a}$ on axis $a$.

    \subsection{Noise-Only Distance for Two Genes}

        Consider two genes $g_1$ and $g_2$ with mean expression vectors
        \[
            \vec{\mu}_1 = (\bar{x}_{g_1,1}, \ldots, \bar{x}_{g_1,A}), \quad
            \vec{\mu}_2 = (\bar{x}_{g_2,1}, \ldots, \bar{x}_{g_2,A}),
        \]
        where $A$ is the number of axes.

        The observed Euclidean distance is
        \[
            D_{\mathrm{obs}} = \left\| \vec{\mu}_1 - \vec{\mu}_2 \right\|.
        \]

        To estimate how much distance can be generated by noise alone, we perform Monte Carlo sampling:

        For each iteration $b$:

        \begin{enumerate}
            \item For each axis $a$, sample independently with replacement:
                  \[
                      \eta_{1,a}^{(b)} \sim \mathcal{E}_{g_1,a}, \quad
                      \eta_{2,a}^{(b)} \sim \mathcal{E}_{g_2,a},
                  \]
                  where $\mathcal{E}_{g,a}$ denotes the empirical noise distribution obtained from the kNN bin for gene $g$ on axis $a$.
            \item Compute the noise-only difference for each axis:
                  \[
                      d_a^{(b)} = \eta_{1,a}^{(b)} - \eta_{2,a}^{(b)}.
                  \]
            \item Assemble the noise-only difference vector
                  \[
                      \vec{d}^{(b)} = (d_1^{(b)}, \ldots, d_A^{(b)}),
                  \]
                  and compute its Euclidean norm:
                  \[
                      D_{\mathrm{null}}^{(b)} = \left\| \vec{d}^{(b)} \right\|.
                  \]
        \end{enumerate}

        Repeating this $B$ times (default $B=100$ or $B=1000$) yields an empirical null distribution $\{D_{\mathrm{null}}^{(b)}\}$ of distances generated purely by replicate variability.

    \subsection{Empirical P-Value}

        The empirical p-value for the observed distance is computed as
        \[
            p = \frac{1 + \sum_{b=1}^{B} I\!\left(D_{\mathrm{null}}^{(b)} \ge D_{\mathrm{obs}}\right)}{B + 1},
        \]
        where $I(\cdot)$ is the indicator function.

        This p-value estimates the probability that a distance at least as large as the observed one could arise from biological and technical replicate variability alone.

        Distances with small $p$ are therefore unlikely to be explainable by noise and indicate robust biological separation.

        \subsubsection{Noise distribution quantile}

            The noise percentile an observed distance falls into is computed as:
            \[
                q = \frac{1 + \sum_{b=1}^{B} I\!\left(D_{\mathrm{null}}^{(b)} \le D_{\mathrm{obs}}\right)}{B + 1}.
            \]

            This value is used in the below plots 3 and 4 (see Sections~\ref{ssec:noise_vs_obs_distance_plot} and~\ref{ssec:obs_vs_noise_quant_plot}).

    \subsection{Diagnostic Plots for Empirical Noise Calibration}

        To ensure that the empirical noise model behaves as intended and that distance significance estimates are trustworthy, we generate three diagnostic plots. These plots are designed to (i) visualize the raw replicate structure, (ii) verify the statistical properties of the sampled noise residuals, and (iii) directly compare observed distances to their noise-only null distributions. All plots are generated per axis (tissue/condition) unless stated otherwise.

        \subsubsection{Plot 1: Mean Expression vs Biological Replicate Values (Scatterplot)}

            \paragraph{Why we plot it.}
            This plot provides a direct view of the raw data structure used by the empirical noise model. It answers the basic question: how do biological replicate values vary as a function of mean expression level on a given axis? In particular, it reveals heteroscedasticity (mean-dependent variance), heavy tails, and potential artifacts such as floor effects at low expression.

            \paragraph{How we generate it.}
            For each axis $a$:
            \begin{enumerate}
                \item For each gene $g$, compute the mean expression $\bar{x}_{g,a}$ across biological replicates.
                \item For each replicate $i$ of gene $g$, emit one point with coordinates:
                      \[
                        (\bar{x}_{g,a}, x_{g,a,i}).
                      \]
            \end{enumerate}

            \paragraph{What the plot shows and informs about.}
            The scatterplot shows the relationship between mean expression and replicate variability on that axis. It informs whether variability increases at low expression (common for sparse or low-abundance genes), whether the spread is symmetric, and whether there are systematic deviations (outliers or saturation effects). This plot is the primary sanity check that the empirical neighborhood sampling is grounded in the actual replicate structure.

        \subsubsection{Plot 2: Mean Expression vs Signed Residuals (Residual Cloud Plot)}

            \paragraph{Why we plot it.}
            The empirical null distribution is sampled from \emph{signed} residuals
            \[
                \epsilon = x_{g,a,i} - \bar{x}_{g,a}.
            \]
            This plot verifies that residuals behave in a biologically plausible way and that sampling signed residuals (rather than absolute or squared deviations) is justified. It also reveals whether residual distributions are skewed, heavy-tailed, multi-modal, or contain systematic artifacts, which directly affects null distance calibration.

            \paragraph{How we generate it.}
            For each axis $a$:
            \begin{enumerate}
                \item For each gene $g$, compute the mean expression $\bar{x}_{g,a}$.
                \item For each replicate $i$, compute the signed residual:
                      \[
                          \epsilon_{g,a,i} = x_{g,a,i} - \bar{x}_{g,a}.
                      \]
                \item Emit one point with coordinates:
                      \[
                          (\bar{x}_{g,a}, \epsilon_{g,a,i}).
                      \]
            \end{enumerate}
            Optionally, for visualization clarity, points may be alpha-blended or downsampled uniformly, but the residual computation must always use the full replicate set.

            \paragraph{What the plot shows and informs about.}
            The residual cloud plot shows how the noise distribution depends on expression level. A symmetric cloud around zero indicates that signed residual sampling will allow both cancellation and amplification when forming noise-only differences. Strong skewness suggests asymmetric variability (biological or technical), while heavy tails indicate that extreme null distances may occur more often than Gaussian assumptions would predict. Multi-modality can indicate batch effects, mixed subpopulations, or inconsistent replicate definitions. In all cases, the plot informs how conservative or liberal the empirical null may become and whether additional stratification (e.g. by project or batch) is required.

        \subsubsection{Plot 3: Observed Distance vs Noise-Only Null Distance Distribution}\label{ssec:noise_vs_obs_distance_plot}

            This plot applies for the assessment of significance of a specific Euclidean Distance e.g. comparing two genes or a family centroid with a particular gene.

            \paragraph{Why we plot it.}
            The final output of the method is an empirical p-value that assesses whether an observed Euclidean distance can be explained by replicate variability alone. This plot provides an intuitive, model-independent visualization of that comparison. It answers: for a given gene pair, is the measured distance clearly separated from the noise-only distances, or does it lie within the typical null range?

            \paragraph{How we generate it.}
            For a fixed pair of genes $(g_1, g_2)$:
            \begin{enumerate}
                \item Compute the observed distance
                      \[
                      D_{\mathrm{obs}} = \left\| \vec{\mu}_1 - \vec{\mu}_2 \right\|,
                      \]
                      with mean expression vectors $\vec{\mu}_1, \vec{\mu}_2$ as defined in Section~\ref{sec:vecspace}.
                \item Generate a null distribution of noise-only distances $\{D_{\mathrm{null}}^{(b)}\}_{b=1}^B$ using the Monte Carlo procedure described in the previous subsection, sampling signed residuals from local kNN bins per axis.
                \item Visualize the null distribution for this gene pair (e.g. as a violin plot, density plot, or box plot) and overlay the observed distance $D_{\mathrm{obs}}$ as a vertical line or point.
            \end{enumerate}
            For screening many gene pairs, the same concept can be summarized by plotting null quantiles (e.g. 50th, 95th, 99th percentiles) against $D_{\mathrm{obs}}$.

            \paragraph{What the plot shows and informs about.}
            This plot shows whether the observed distance is large relative to the noise-only null distribution. If $D_{\mathrm{obs}}$ lies far above the upper tail of the null, the distance is robust and unlikely to be explainable by replicate variability alone. If $D_{\mathrm{obs}}$ overlaps the bulk of the null, the separation is not robust given the replicate structure and should be interpreted cautiously. This plot also reveals cases where replicate imbalance (or mean-dependent variability) leads to broad null distributions, making it clear why p-values become large even when raw distances look substantial.

        \subsubsection{Plot 4: Observed Distance vs Noise Percentile (Global Scatterplot)}\label{ssec:obs_vs_noise_quant_plot}

            This is a generalization of Plot 3 (see Section~\ref{ssec:noise_vs_obs_distance_plot}).

            \paragraph{Why we plot it.}
            The empirical noise model yields, for each observed Euclidean distance, a null distribution of noise-only distances. While this can be visualized for a single gene pair (e.g. as a raincloud or violin plot), it is often more informative to assess the behavior of the noise model across many observed distances. This plot provides a global diagnostic that answers: as observed distances increase, do they systematically move into the upper tail of their corresponding noise-only null distributions, as expected for true signal? It also reveals whether many observed distances are still explainable by noise, which would indicate insufficient replication or high biological variability.

            \paragraph{How we generate it.}
            \begin{enumerate}
                \item During the calculation of all Euclidean Distances e.g. while iterating over all gene families, store all observed distances ($D_{\mathrm{obs}}$) and the percentile of the Null noise model ($q$) in which the observed distances fell.
                \item Plot the points $(D_{\mathrm{obs}}, q)$ as a scatterplot, where the x-axis shows the observed distance and the y-axis shows the noise percentile $q$.
            \end{enumerate}

            \emph{Note} if the number of Euclidean distances is very large, a random subset of observed distances can be stored for visualization.

            \paragraph{What the plot shows and informs about.}
            Each point summarizes one observed expression distance, either from a pair of genes or from a family centroid to gene: $D_{\mathrm{obs}}$ is the measured separation in semantic expression space, and $q$ indicates where this distance falls within the noise-only null distribution for that specific pair. Values near $q \approx 0.5$ indicate that the observed distance is typical under noise alone, while values near $q \approx 1$ indicate that the observed distance lies in the extreme upper tail of the null and is therefore unlikely to be explainable by replicate variability.

            The overall trend of this scatterplot provides a global calibration check. In well-calibrated settings, larger observed distances should increasingly concentrate at high percentiles. A substantial mass of large distances with only moderate percentiles suggests that replicate variability (biological or technical) is large relative to the measured separation, which may arise from low replication on some axes or high intrinsic inter-individual variability in the sampled conditions.